\section{Supplementary methods and reproducibility}
All supplementary statistics use the corrected unloaded-subject reference $\mathbf X+\boldsymbol\phi_i$ and loaded coordinates $\mathbf X+\boldsymbol\phi_i+\mathbf u_{i\ell}$, as specified in the mathematical appendix of the main article. The earlier template-to-subject extraction is superseded. The three simulation variants share subject indices; resampling preserves their pairing.

Component differences compare bilateral absolute magnitudes, not signed motion vectors. Their bootstrap intervals describe medians; primary tests are exact binomial sign tests. Signed-rank results are retained only as machine-readable sensitivity diagnostics. The six fully adjusted sex comparisons use HC3 covariance and a single six-test FDR family. Variance ratios are non-additive and are not causal contributions.

The scaling model is
\begin{align}
\log Y&=\beta_0+\beta_v\log V+\beta_s\mathrm{Female}+\beta_{vs}\log V\,\mathrm{Female}\nonumber\\
&+\beta_a\mathrm{Age}_z+\beta_{AP}\mathrm{AP}_z+\beta_b\mathrm{Biischiadic}_z+\beta_p\mathrm{Subpubic}_z+\epsilon.
\end{align}
Male and female slopes are $\beta_v$ and $\beta_v+\beta_{vs}$. Allometry slope tests are corrected jointly across 20 sex-specific slopes; the ten interaction tests form a separate family. Table S4 is an exact reparameterization with $\log s=(\log V-\log V_{template})/3$, so scale exponents equal three times the volume exponents. It is not an independent validation or sensitivity analysis.

The accompanying reproduction guide specifies the complete command sequence for corrected extraction, extraction sensitivity, statistics, table formatting, numerical manuscript passages and publication figures. Original FE displacement archives are unchanged. Corrected checkpoints are stored separately from statistical outputs. Existing checkpoints must be removed or regenerated if the source fields or extraction definitions change.
